\documentclass[10pt]{article}

\usepackage[T1]{fontenc}
\usepackage{amsmath,amssymb,slashed}
\usepackage[margin=0.75in]{geometry}

\newcommand{\Tr}{\operatorname{Tr}}
\newcommand{\id}{\mathbb{I}}

\begin{document}
\begin{center}
{\Large\bfseries Spin-Density Matrix and Fragmentation Cross Section}
\end{center}
\vspace{0.25em}
\small

\section{Hard spin density matrix and Fragmentation Cross Section}

In a common Cartesian spin basis,
\[
\rho=\frac14\left[
\id\otimes\id+B_i\sigma^i\otimes\id
+\bar B_j\id\otimes\sigma^j
+C_{ij}\sigma^i\otimes\sigma^j\right],
\qquad i,j=x,y,z .
\]
Thus \(B_i\), \(\bar B_j\), and \(C_{ij}\) are respectively the quark
polarization, antiquark polarization, and spin-correlation tensor. Insert the two fragmentation spin factors directly on the two spin lines:
\begin{align*}
{\cal K}
&=\Tr_{s\bar s}\!\left\{\rho\,
(D_1\id+a_i\sigma^i)\otimes
(\bar D_1\id+\bar a_j\sigma^j)\right\}
\\
&=D_1\bar D_1+B_i\bar D_1a_i
+\bar B_jD_1\bar a_j+C_{ij}a_i\bar a_j .
\end{align*}
Here \(\Tr\sigma^i=0\) and
\(\Tr(\sigma^i\sigma^j)=2\delta^{ij}\).
Equivalently,
\[
{\cal K}=w_0+B_iw_{B_i}+\bar B_jw_{\bar B_j}+C_{ij}w_{C_{ij}},
\qquad
(w_0,w_{B_i},w_{\bar B_j},w_{C_{ij}})
=(D_1\bar D_1,\bar D_1a_i,D_1\bar a_j,a_i\bar a_j).
\]

\section{Calculate \(D_1\) and \(a_i\) with the spin projector}

The fragmentation correlator has open Dirac indices:
\begin{align*}
\Delta_{ab}(k;P_1,P_2)=\sum_X\int\frac{d^4\xi}{(2\pi)^4}e^{ik\cdot\xi}
&\langle0|{\cal W}\psi_a(\xi)|P_1P_2,X\rangle
\langle P_1P_2,X|\bar\psi_b(0){\cal W}|0\rangle ,
\end{align*}
evaluated with \(k^+\equiv Q\) fixed (the frame with \(P_{h\perp}=0\),
\(Q\) the large light-cone momentum of the fragmenting quark) and
integrated over the small component \(k^-\).  At leading twist \(\Delta\)
is built entirely from the structures \(\gamma^-\), \(\gamma^-\gamma_5\),
\(i\sigma^{i-}\gamma_5\) (\(i=x,y\)).  \(D_1\) and \(a_i\) are the
coefficients of these structures, read off by contracting \(\Delta\) with
the \emph{conjugate} light-cone direction \(n\) (\(n^2=0\), \(n\cdot k=1\),
\(\slashed n=\gamma^+/Q\)): since \(\Tr[\gamma^+\gamma^-]\neq0\) while
\(\Tr[\gamma^-\gamma^-]=0\), only \(\gamma^+\)-type insertions have a
nonzero overlap with \(\Delta\).  Formally, write
\[
{\cal D}(S)=\frac1{4z}\int dk^-\Tr_D\!\left[\Delta\,\gamma^+(1+\gamma_5\slashed S)\right]
\Big|_{\text{transverse }S}
=D_1+S_ra_r ,
\qquad r=x,y ,
\]
which, using the exact identity \(\gamma^+\gamma_5\slashed e_r=i\sigma^{r+}\gamma_5\)
(\(\slashed e_r=-\gamma^r\)), gives
\[
D_1=\frac1{4z}\int dk^-\Tr_D[\gamma^+\Delta],
\qquad
a_r=\frac1{4z}\int dk^-\Tr_D[i\sigma^{r+}\gamma_5\Delta],
\quad r=x,y .
\]
The longitudinal component is probed directly by \(\gamma_5\) alone,
without an extra slashed unit vector (inserting \(\gamma_5\slashed e_z\)
would produce a structure orthogonal to every term in \(\Delta\)):
\[
a_z=\frac1{4z}\int dk^-\Tr_D[\gamma^+\gamma_5\Delta] .
\]
Evaluating these three traces against the leading-twist parametrization
of \(\Delta\) gives
\begin{align*}
D_1&=\frac1{4z}\int dk^-\Tr_D[\gamma^+\Delta],
\\
a_z&=\frac1{4z}\int dk^-\Tr_D[\gamma^+\gamma_5\Delta]
=\frac{\epsilon_T^{mn}R_{Tm}k_{Tn}}{M_1M_2}G_1^\perp,
\\
a_r&=\frac1{4z}\int dk^-\Tr_D[i\sigma^{r+}\gamma_5\Delta]
=\frac{\epsilon_T^{rn}R_{Tn}}{M_1+M_2}H_1^{\sphericalangle}
+\frac{\epsilon_T^{rn}k_{Tn}}{M_1+M_2}H_1^\perp,
\quad r=x,y .
\end{align*}

For the reduced choice \(H_1^\perp=0\), with
\(\boldsymbol k_T=-\boldsymbol P_{hT}/z\) and
\(\epsilon_T^{xy}=1\), abbreviate
\[
\kappa=\frac{R_T}{M_1+M_2}H_1^{\sphericalangle},
\qquad
\lambda=\frac{R_TP_{hT}}{zM_1M_2}G_1^\perp,
\]
so that
\[
\mathbf a\equiv\begin{pmatrix}a_x\\a_y\\a_z\end{pmatrix}
=\begin{pmatrix}\kappa\sin\phi_R\\-\kappa\cos\phi_R\\
\lambda\sin(\phi_R-\phi_h)\end{pmatrix} .
\]
\paragraph{Antiquark side.} The antiquark moves in the \emph{opposite}
direction, with large light-cone component \(k'^-\equiv\bar Q\) instead
of \(k'^+\).  Its correlator \(\bar\Delta\) is therefore built at leading
twist from the \emph{mirror} structures \(\gamma^+,\gamma^+\gamma_5,
i\sigma^{i+}\gamma_5\), and \(\bar D_1,\bar a_i\) must be extracted with
the conjugate projectors \(\gamma^-,\gamma^-\gamma_5,i\sigma^{r-}\gamma_5\)
— every \(+\leftrightarrow-\) exchanged relative to the quark case:
\[
\bar D_1=\frac1{4\bar z}\int dk'^+\Tr_D[\gamma^-\bar\Delta],
\quad
\bar a_z=\frac1{4\bar z}\int dk'^+\Tr_D[\gamma^-\gamma_5\bar\Delta],
\quad
\bar a_r=\frac1{4\bar z}\int dk'^+\Tr_D[i\sigma^{r-}\gamma_5\bar\Delta].
\]
Charge conjugation relates the antiquark correlator to the quark one with
an extra relative minus sign on every spin-dependent (\(G_1^\perp\)- and
\(H_1\)-type) structure, while the spin-independent \(D_1\) term is
unaffected.  Hence barring every fragmentation variable is \emph{not}
sufficient:
\[
\bar D_1=D_1\big|_{\rm barred},
\qquad
\bar a_i=-\,a_i\big|_{\rm barred} ,
\qquad i=x,y,z ,
\]
where "barred" means every kinematic variable and fragmentation function
on the right (\(z,M_1,M_2,R_T,P_{hT},\phi_R,\phi_h,H_1^{\sphericalangle},
G_1^\perp\)) is replaced by its antiquark counterpart.  With
\(\bar\kappa,\bar\lambda\) denoting \(\kappa,\lambda\) barred, this gives
\[
\bar{\mathbf a}\equiv\begin{pmatrix}\bar a_x\\\bar a_y\\\bar a_z\end{pmatrix}
=\begin{pmatrix}-\bar\kappa\sin\bar\phi_R\\\bar\kappa\cos\bar\phi_R\\
-\bar\lambda\sin(\bar\phi_R-\bar\phi_h)\end{pmatrix} .
\]
If only one fragmentation is observed, the spin factor reduces to
\({\cal K}_{\rm one}=D_1+B_i a_i\).

\section{Spin weights in terms of the fragmentation functions}

With \(\mathbf a,\bar{\mathbf a}\) as defined in Sec.~2, insert the
results of Sec.~2 (with \(H_1^\perp=0\)) into
\(\mathbf w_B=\bar D_1\mathbf a\), \(\mathbf w_{\bar B}=D_1\bar{\mathbf a}\),
and \(W_C=\mathbf a\,\bar{\mathbf a}^{\,T}\).  The single-spin weight
vectors are
\[
\mathbf w_B=\bar D_1
\begin{pmatrix}\kappa\sin\phi_R\\-\kappa\cos\phi_R\\
\lambda\sin(\phi_R-\phi_h)\end{pmatrix},
\qquad
\mathbf w_{\bar B}=D_1
\begin{pmatrix}-\bar\kappa\sin\bar\phi_R\\\bar\kappa\cos\bar\phi_R\\
-\bar\lambda\sin(\bar\phi_R-\bar\phi_h)\end{pmatrix},
\]
and the correlation-weight matrix is the outer product
\[
W_C=\mathbf a\,\bar{\mathbf a}^{\,T}=
\begin{pmatrix}
-\kappa\bar\kappa\sin\phi_R\sin\bar\phi_R &
\kappa\bar\kappa\sin\phi_R\cos\bar\phi_R &
-\kappa\bar\lambda\sin\phi_R\sin(\bar\phi_R-\bar\phi_h)
\\[2pt]
\kappa\bar\kappa\cos\phi_R\sin\bar\phi_R &
-\kappa\bar\kappa\cos\phi_R\cos\bar\phi_R &
\kappa\bar\lambda\cos\phi_R\sin(\bar\phi_R-\bar\phi_h)
\\[2pt]
-\lambda\bar\kappa\sin(\phi_R-\phi_h)\sin\bar\phi_R &
\lambda\bar\kappa\sin(\phi_R-\phi_h)\cos\bar\phi_R &
-\lambda\bar\lambda\sin(\phi_R-\phi_h)\sin(\bar\phi_R-\bar\phi_h)
\end{pmatrix} .
\]
Every entry of \(\mathbf w_B\), \(\mathbf w_{\bar B}\), and \(W_C\) is thus
expressed directly in terms of the collinear-times-TMD fragmentation
functions \(D_1\), \(H_1^{\sphericalangle}\), \(G_1^\perp\) (and their
barred counterparts) together with the observed hadron-pair kinematics
\(R_T,P_{hT},\phi_R,\phi_h\) (and barred).

\section{Deriving \({\cal K}\) with the quark mass kept}

Specialize Sec.~1 to \(e^+e^-\to\gamma^*\to q\bar q\), with \(\theta\)
the quark production angle and
\[
\rho_m\equiv\frac{4m^2}{s},
\quad
A_0=1+\cos^2\theta,
\quad
B_0=\sin^2\theta,
\quad
A_U=A_0+\rho_mB_0 .
\]
Mass can enter through the hard vertex (\(B_i,\bar B_j,C_{ij}\)) or the
fragmentation vertex (\(D_1,a_i\)).

Contracting the leptonic tensor with the covariant spin sums
\(u(k,S)\bar u(k,S)=\tfrac12(\slashed k+m)(1+\gamma_5\slashed S)\),
\(v(k',S')\bar v(k',S')=\tfrac12(\slashed k'-m)(1+\gamma_5\slashed S')\)
at the photon vertex gives the standard exact-in-\(m\) result: parity
still forces \(B_i=\bar B_j=0\), and in the \((x,y,z)=(-r,n,k)\) axes of
Sec.~1 (\(k\) along the quark momentum, \(r\) in the lepton-quark plane,
\(n\) normal to it),
\[
C=\frac1{A_U}\begin{pmatrix}
(1+\rho_m)B_0 & 0 & -2\sqrt{\rho_m}\sin\theta\cos\theta\\
0 & -(1-\rho_m)B_0 & 0\\
-2\sqrt{\rho_m}\sin\theta\cos\theta & 0 & A_0-\rho_mB_0
\end{pmatrix} .
\]

\subsection*{Mass terms from the final-state quark propagator}

Decompose the exact final-state spin sum by light-cone power,
\[
\slashed k+m
=k^+\gamma^-+(\slashed k_\perp+m)+k^-\gamma^+,
\qquad
k^-=\frac{m^2+\boldsymbol k_\perp^2}{2k^+} ,
\]
and dress \(\Delta\)'s known structures with the full propagator instead
of its truncation \(k^+\gamma^-\):
\[
\Delta_{\rm massive}=(\slashed k+m)D_1
+\gamma_5\gamma^0\gamma^3(\slashed k+m)\,X_zG_1^\perp
+\gamma^r\gamma_5(\slashed k+m)Y_r ,
\]
with \(X_z,Y_r\) as in Sec.~2.  Every trace of the subleading
\((\slashed k_\perp+m)\) and \(k^-\gamma^+\) blocks against the Sec.~2
projectors \(\gamma^+,\gamma^+\gamma_5,i\sigma^{r+}\gamma_5\) vanishes
identically (verified by explicit Dirac-trace computation, keeping
\(m,k^-,\boldsymbol k_\perp\) fully general):
\[
\Tr_D[\gamma^+(\slashed k_\perp+m)]=\Tr_D[\gamma^+\gamma^+]
=\Tr_D[\gamma^+\gamma^0\gamma^3(\slashed k_\perp+m)]
=\Tr_D[\gamma^+\gamma^0\gamma^3\gamma^+]
=\Tr_D[i\sigma^{r+}\gamma_5(\slashed k_\perp+m)]
=\Tr_D[i\sigma^{r+}\gamma_5\gamma^+]=0 .
\]
So \(D_1,a_z,a_r\) are unchanged at any order in \(m\) or
\(\boldsymbol k_\perp\): the quark mass cannot enter through these three
Dirac channels at all.

It can only enter through a channel Sec.~2 never uses: the identity.
The scalar term in \(\slashed k_\perp+m\) generates
\(\Delta\supset\mathbb I\cdot E(z,\boldsymbol k_\perp;m)\), read off by
the scalar trace \(E=\tfrac1{4z}\int dk^-\Tr_D[\Delta]\).  Since
\(\Tr_D[\slashed k+m]=0+4m\), this ansatz gives \(E=mD_1\) directly from
the \(+m\) in the propagator.  But
\({\cal K}=\Tr_{s\bar s}[\rho(D_1\mathbb I+a_i\sigma^i)\otimes
(\bar D_1\mathbb I+\bar a_j\sigma^j)]\) traces over quark-spin space, not
Dirac space, so \(E\) cannot correct \(D_1\) or \(a_i\) inside it — it
would enter only through a separate \(E\bar D_1+D_1\bar E\) term
alongside \({\cal K}\), left out below.

\subsection*{Result}

\(D_1,a_i,\bar D_1,\bar a_i\) are mass-independent, so all of
\({\cal K}\)'s mass dependence sits in \(C\).  With \(B_i=\bar B_j=0\),
\({\cal K}=D_1\bar D_1+C_{ij}a_i\bar a_j\) is the bilinear form
\[
{\cal K}=D_1\bar D_1+\mathbf a^T C\,\bar{\mathbf a} ,
\]
using the same \(\mathbf a,\bar{\mathbf a}\) of Sec.~2.  Expanding the
product gives
\[
\begin{aligned}
{\cal K}={}&D_1\bar D_1
+\frac1{A_U}\Big[
B_0\kappa\bar\kappa\cos(\phi_R+\bar\phi_R)
-A_0\lambda\bar\lambda\sin(\phi_R-\phi_h)\sin(\bar\phi_R-\bar\phi_h)
\Big]
\\[2pt]
&+\frac{\sqrt{\rho_m}\sin2\theta}{A_U}\Big[
\kappa\bar\lambda\sin\phi_R\sin(\bar\phi_R-\bar\phi_h)
+\lambda\bar\kappa\sin(\phi_R-\phi_h)\sin\bar\phi_R
\Big]
\\[2pt]
&+\frac{\rho_mB_0}{A_U}\Big[
\lambda\bar\lambda\sin(\phi_R-\phi_h)\sin(\bar\phi_R-\bar\phi_h)
-\kappa\bar\kappa\cos(\phi_R-\bar\phi_R)
\Big] ,
\end{aligned}
\]
whose three lines are the massless Collins-type correlation, an
\(O(\sqrt{\rho_m})\) term mixing the transverse sector of one jet with
the longitudinal sector of the other (vanishing at \(m=0\)), and an
\(O(\rho_m)\) correction with the same sum-to-difference angle flip
found in the exact massive double-transverse and double-Collins
\(e^+e^-\) cross sections.

\end{document}
